\section{The crystalline lattice and the exponent on a split column}

% BEGIN S79P1
\paragraph{S79P1. The local field and lattice at a declared column.}
Use the hub, parameter, comparison data and $\mathrm{AdmJ}(L)$
of F78D2. Fix a twisted localization $i$ with split prime
$\ell=\ell_i$ outside the two distinguished primes. Suppose $u$
belongs to the declared column at $i$, and $\pi\in Z_{\rm good}(i)$.
This includes every split column in C78B5. We give the data for
these very $u,\pi$.

DFD8 identifies its declaration with $\mathrm{dcOf}$. The default
declaration has no columns, so the witnessed split declaration was
used. Its covering identity puts $u$ over $\ell$. Its split guard
gives $\ell\ge5$ and $\ell\ne q,d$. The split class is disjoint
from the host's small-prime class. Hence $\ell$ does not divide
the host conductor $N$. Take the original bridge $K=K_u$ and
the bridge coefficient field $L_u$ itself. The prescribed embedding
of $E$ has image in the unramified extension of $\mathbb Q_\ell$:
the host realization is generated inside the cyclotomic field of
conductor $N$, and the $N$th roots of unity are unramified at $\ell$.
The intersection of $K$ with the maximal unramified extension is its
finite unramified subfield, closed as a finite-dimensional
$\mathbb Q_\ell$-subspace. It contains the dense image of $E$,
so it equals $K$. Likewise the $\mathbb Q_\ell$-span of the images
of a $\mathbb Q$-basis of $E$ is closed and contains $K$. Thus
$[K:\mathbb Q_\ell]\le[E:\mathbb Q]=4n$. The bridge's splitting
condition on $L_u$, $\ell\ne2$, and $K=K_0\subset K_0(\zeta_\ell)$
give its good layer at level one. In particular $L_u$ contains
the images of all embeddings of $K$.

Write $\mathcal O=\mathcal O_{L_u}$. Take the standard lattice
$T=\mathcal O^r$ inside $L_u^r$. Its action is obtained by
restricting $\chi_1^{-1}\rho_i(\pi)$ to the decomposition group at
$u$ and applying the prescribed coefficient map to each entry.
The original $\rho_i(\pi)$ and $\chi_1$ are continuous and
integral, with the character unit-valued. Hence these matrices and
their inverses preserve $T$, and the resulting action is
continuous. This is exactly $\mathrm{wLat}$ for the self coefficient
extension. Its action and its lattice give the coefficient bridge of
C78B5 on passage to the quotient $L_u^r/T$.
% END S79P1

% BEGIN S79P2
\paragraph{S79P2. Crystallinity at the selected good point.}
Write $R_\pi$ for the $n_V$-dimensional point representation.
The canonically chosen point family of DFD7--DFD9 gives its
$\mathrm{GoodR}$ property and the complementary representation
$\rho_i(\pi)$ at every $\pi\in Z(i)$. The good-point tie puts our
$\pi$ in $Z(i)$. On fraction spaces its splitting has the form
\[
 R_\pi \simeq \rho_i(\pi)\oplus
                       \bigoplus_{c=1}^{t}\chi_c .
\]
Here the characters are indexed as in the source, including $\chi_1$.
Indeed the $\rho$ field of the complement condition gives an
invertible linear change of coordinates intertwining each original
matrix with the action on this direct sum.
The inclusion and retraction of each block give intertwining
matrices $B,A$ with $AB=1$. Restricting to $G_K$ and mapping the
matrix entries into $L_u$ preserves these identities.

Away from the two distinguished primes, $\mathrm{Conclusion41}$
on this $\mathrm{GoodR}$ family gives crystallinity of $R_\pi$
at $u$. This is the first constituent of $\mathrm{Conclusion1}$
in the setup facts. The local crystalline tie applies to this
continuous matrix on the self coefficient bridge of S79P1.
The two summands $\rho_i(\pi)$ and $\chi_1$ therefore are
crystalline. Over the tied crystalline period ring one can check
this directly on the splitting: the period invariants of a direct
sum are the direct sum of the invariants, and the period dimensions
are bounded above by the corresponding ranks. Equality for the
sum forces equality for each summand. Equivalently the crystalline
comparison isomorphism restricts to the two block idempotents.
Taking the dual of the rank-one crystalline character and then
tensoring with the crystalline $\rho_i(\pi)$ preserves
crystallinity, by dualizing and tensoring the period bases.
This proves crystallinity for the action on $T$ in S79P1, with
the inverse twist in $\mathrm{WpO}$.
% END S79P2

% BEGIN S79P3
\paragraph{S79P3. The jump multiset after removing the character block.}
Take any $\tau:K\to L_u$ over $\mathbb Q_\ell$ and set $j=j_i$.
The $\mathrm{AdmRead}$ and $\mathrm{AdmDict}$ fields of the supplied
$\mathrm{AdmJ}(L)$ choose the complex embedding above $\tau$, its
archimedean place, a sign $\sigma=\pm1$ and an integer $c_\tau$.
Indeed a host embedding lies in the chosen CM half or its conjugate,
which supplies the place used by the dictionary. At this place
write $s=e_{j-1}$, $d_0=(t-1)/2$, and $v_a,b_a$ for the original
archimedean parameters. The untwisted point representation has
the ambient Hodge multiset, since $\ell$ is away from the two
distinguished primes, $\mathrm{GoodR}$ holds and the
de Rham/Weil--Deligne tie supplies the de Rham reading. The dictionary
lists this multiset as $\sigma p+c_\tau$, for the multiset of entries
\[
 p=v_a\ (0\le a<k-1),\qquad p=-b_a\ (0\le a<m),\qquad
 p=-s+d_0-c\ (0\le c<t).
\]
The same dictionary assigns the single weight
$c_\chi=-\sigma s-d_0+c_\tau$ to $\chi_1$. The reading of the
ambient inverse twist at $a$ is its old multiplicity at $a+c_\chi$.
Thus the twisted weights are $d_0+\sigma(p+s)$. For positive sign
they are
\[
 \{s+v_a+d_0\}_a,\qquad \{s-b_a+d_0\}_a,\qquad
 \{0,\ldots,t-1\}.
\]
For negative sign reflect each entry $w$ about $d_0$, sending it to
$2d_0-w$. This permutation preserves the third block. In the
semisimple decomposition of S79P2 the $\mathrm{AdmRead}$ identity
subtracts one from the multiplicity at each integer $0\le a<t$,
and subtracts zero elsewhere, to give the multiplicity for
$\mathrm{WpO}$. The remaining multiset therefore consists of
exactly the first two lists, with the indicated reflection where
required. The jump criterion in that same reading is nonzero
multiplicity. This computation needs only the displayed fields of
$\mathrm{AdmJ}(L)$ and the point splitting.
% END S79P3

% BEGIN S79P4
\paragraph{S79P4. The gap at zero and a uniform height.}
Put $b=\exp(x^\epsilon)$, $\eta=L.\mathrm{etaSz}$, $S=L.\mathrm{Sc}$.
The archimedean size and gap inequalities on the place of S79P3 give
\[
 \begin{gathered}
 0\le v_a\le S,\qquad
 (t+3)/2\le |s-b_a|\le\eta,\qquad
 \eta+t+S<s\le b+\eta+1,\\
 \eta+t+S\le b .
 \end{gathered}
\]
For the low weights $s=e_{j-1}$, and the ordering of $b_a$ together
with the two neighboring $b,e$ gaps gives the lower bound on either
side of $j-1$. The upper bound is the size datum's
$|e_{j-1}-b_a|\le\eta$. For the bound on $s$ use the thick, top
and thin windows for $j$. They all give $s\ge w''$ and
$s\le b+\eta+1$: the top window uses
$b_\nu+b'/x\le b$ from $\mathrm{AdmPar}$, the thin one uses
$3w''\le b$, and the thick window has center $b$. Finally
$\eta+t+S<w''$ and $3w''\le b$ are part of that same admissibility.

The first two lists of S79P3 have entries and reflected entries of
the form $d_0\pm v$ with $v=s+v_a$ or $v=s-b_a$.
In the first case $v>t$, and in the second
$|v|\ge(t+3)/2=d_0+2$. Hence zero is absent for both signs.
Their absolute values are at most
$b+\eta+1+S+d_0\le2b$, because $1+d_0\le t$.
The number of entries for each embedding is at most
$(k-1)+m=r$, even if two entries coincide. Retain their integral
values in a finite set $U_\tau$. All jumps lie there. For
$h=\lfloor2b\rfloor+1$ they lie in $[-h,h]$, and in fact in
$[-(h-1),h-1]$. Since $K$ splits in $L_u$, its tensor with $L_u$
is the product of the factors for these embeddings. The idempotents
split the filtered de Rham space and preserve every filtration
step. A change in the total filtration therefore changes some
embedding component. Thus the total jump set is in
$\bigcup_\tau U_\tau$, of size at most $[K:\mathbb Q_\ell]r\le4nr$.
The same $h$ works for all the split places and good points of
this hub.
% END S79P4

% BEGIN S79P5
\paragraph{S79P5. Comparing the geometric exponent with the stored column exponent.}
Write $J=\lfloor\log_\ell(2h)\rfloor$ for the natural-number
logarithm, and let $n_0$ be the least integer $v\ge0$ such that
$2h\le\ell^v(\ell-1)$. Since $2h<\ell^{J+1}$, $n_0\le J+1$.
The source's cyclotomic level is $m_0=2+n_0+1\le J+4$.
Every nonzero jump $a$ has
$v_\ell(a)\le\lfloor\log_\ell h\rfloor\le J$. Thus for each
embedding the actual geometric exponent has the bound
\[
 c_0^{\rm loc}(\tau)
    =2+\sum_{a\text{ a jump at }\tau}\bigl(m_0+v_\ell(a)\bigr)
    \le 2+r(4+2J).
\]
Here the jumps are counted once per distinct integer. S79P4 removes
zero before this valuation estimate.

One also obtains the source's bound using its total jump set.
For the unramified $K$ the fields
$K_0(\zeta_{\ell^v})$ for $v\ge1$ have degrees
$[K:\mathbb Q_\ell](\ell-1)\ell^{v-1}$ over $\mathbb Q_\ell$.
Indeed $\Phi_{\ell^v}(1+X)$ is Eisenstein and remains so over
the unramified $K$. The difference of the $\ell$-adic valuations
of the degree at $m_0$ and of $[K:\mathbb Q_\ell]$ is $m_0-1$.
In the definition of $\mathrm{cLL}$ the outer sum over $[0,1)$
has only index zero. By S79P4 its inner sum is bounded by
$4nr(4+2J)$. The maximum of this bound and one equals this bound
because $n,r\ge1$. Hence
\[
  \mathrm{invariantsCountConstant}
    \le (m_0-1)+2+4nr(4+2J)
    \le 5+J+4n n_V(4+2J)=X.c_0(i,u).
\]
The last identity is DFD8 on this very hub and $r\le n_V$.
It also bounds $c_0^{\rm loc}(\tau)$. Together S79P1--S79P5 give
every field of the split-column witness
$\mathrm{JointCond\_H0Wall}$ with $M_0=h_0=h$.
Consequently any local theorem killing $(L_u^r/T)^{G_K}$ by
$\ell^{c_0^{\rm loc}(\tau)}$, or by the collective exponent just
bounded, supplies the exponent required in C78B5 by multiplication
by the remaining power of $\ell$. The numerical comparison does
not assume an action of this exponent on invariant classes.
% END S79P5
